\clearpage
\appendix
\section{More Details}
\label{Training}
Following the data and reward configurations of Flow-GRPO, we conducted multi-task hybrid training for GRPO-mix using an epoch ratio of 3:1:1 for GenEval, OCR, and PickScore, respectively. During each epoch, rewards were exclusively provided by the reward model corresponding to the current data partition. Training was executed on a distributed cluster of four nodes, each equipped with eight H800 GPUs for about 50 hours. For the GenEval, OCR, and PickScore teachers, we utilized the official Flow-GRPO checkpoints. Additionally, to incorporate the DeQA teacher—which focuses solely on image quality—its reward signals were integrated into the standard GRPO-mix training via a 1:1 summation ratio.

\paragraph{Hyperparameters Specification}\label{app:subsubsec:synthetic-hyperp}
Except for $\beta$, GRPO hyperparameters are fixed across tasks. We use a sampling timestep $T=10$ and an evaluation timestep $T=40$. Other settings include a group size $G=24$, an noise level $a=0.7$ and an image resolution of 512. The MAR KL ratio $\beta$ is set to 0.02. We use Lora with $\alpha=64$ and $r=32$.

\begin{table}[h]
\centering
\resizebox{\linewidth}{!}{
\begin{tabular}{ll}
\toprule
\textbf{Models} & \textbf{Links} \\
\midrule
\texttt{SD3.5-M}~\citep{esser2024scaling} & \url{https://huggingface.co/stabilityai/stable-diffusion-3.5-medium} \\
\texttt{Aesthetic Score}~\citep{Aesthetics} & \url{https://github.com/LAION-AI/aesthetic-predictor} \\
\texttt{PickScore}~\citep{pickscore} & \url{https://huggingface.co/yuvalkirstain/PickScore_v1} \\
\texttt{DeQA score}~\citep{deqa} & \url{https://huggingface.co/zhiyuanyou/DeQA-Score-Mix3} \\
\texttt{ImageReward}~\citep{xu2024imagereward} & \url{https://huggingface.co/THUDM/ImageReward} \\
\texttt{UnifiedReward}~\citep{wang2025unified} & \url{https://huggingface.co/CodeGoat24/UnifiedReward-7b-v1.5} \\
\texttt{HPS-v2.1}~\citep{hps} & \url{https://github.com/tgxs002/HPSv2} \\
\texttt{Qwenvl Score}~\citep{yang2025qwen3} & \url{https://huggingface.co/Qwen/Qwen3-30B-A3B-Instruct-2507} \\
\bottomrule
\end{tabular}
}
\end{table}
Regarding Qwenvl Score, we adapt the prompt used in Flow-GRPO~\cite{liu2025flow}. The prompt is shown in Fig.~\ref{fig:evaluation_prompt}. We use 
Qwen3-30B-A3B-Instruct-2507.

\input{tables/eval_prompt}
\section{More Results}
\subsection{Qualitative results}
More qualitative results are shown in Fig.~\ref{fig: result_pick}, ~\ref{fig: result_geneval} and ~\ref{fig: result_ocr}. Our approach not only ensures precise content generation but also delivers superior image quality and coherent structural layouts. By achieving stronger alignment with human preferences, Flow-OPD demonstrates significant novelty in bridging functional accuracy with aesthetic excellence.
\begin{figure*}[!h]
    \centering
    \includegraphics[scale=0.6]{images/more.pdf}
    % \vspace{-15pt}
    \caption{More quantitative comparisons on the Pickscore evaluation set.
    }
    \label{fig: result_pick}
    % \vspace{-0.4cm}
\end{figure*}

\begin{figure*}[!h]
    \centering
    \includegraphics[scale=0.6]{images/more2.pdf}
    % \vspace{-15pt}
    \caption{More quantitative comparisons on the GenEval evaluation set.
    }
    \label{fig: result_geneval}
    % \vspace{-0.4cm}
\end{figure*}


\begin{figure*}[!h]
    \centering
    \includegraphics[scale=0.6]{images/more3.pdf}
    % \vspace{-15pt}
    \caption{More quantitative comparisons on the OCR evaluation set.
    }
    \label{fig: result_ocr}
    % \vspace{-0.4cm}
\end{figure*}
\subsection{Comparison with DiffusionNFT}
DiffusionNFT~\cite{zheng2025diffusionnft} introduces an online reinforcement learning framework that directly integrates reward feedback into the forward diffusion process, enabling effective policy optimization during the noise-injection phase.
Despite achieving competitive benchmark scores, DiffusionNFT exhibits several critical limitations. First, it is fundamentally incompatible with Classifier-Free Guidance (CFG), which severely bottlenecks its performance upper bound. Second, it suffers from pronounced reward hacking. As illustrated in Fig.~\ref{fig: diffusion}
, while the model correctly generates the targeted text and 'sunset' elements, it simultaneously hallucinates malformed hands and extraneous objects (e.g., oranges), accompanied by severe \textit{over-smoothed, plastic-like textural artifacts}. Current standard benchmarks largely overlook these localized structural and aesthetic failures. To address this evaluation blind spot, we employ the Qwenvl Score for a more comprehensive assessment. By leveraging continuous-time dense multi-expert supervision and task-style decoupling, Flow-OPD effectively circumvents these reward hacking behaviors, achieving significantly higher Qwen-VL scores than DiffusionNFT. These findings also underscore a pressing need within the community to develop more robust, fine-grained evaluation paradigms for text-to-image generation.
\begin{table}[h]
    \small
    \centering
    \caption{Comparison of Human Preference Alignment. Our Flow-OPD consistently achieves superior scores in complex visual reasoning and layout coherence, as evaluated by Qwen-VL.}
    \label{tab:reward_metrics_short}
    \renewcommand{\arraystretch}{1.1} 
    \begin{tabular}{lc} 
        \toprule
        \textbf{Model} & \textbf{Qwen-VL Score} \\
        \midrule
        DiffusionNFT & 3.74 \\
        \textbf{Ours (Flow-OPD)} & \textbf{4.05} \\
        \bottomrule
    \end{tabular}
\end{table}
\begin{figure*}[!h]
    \centering
    \includegraphics[scale=0.445]{images/diffusionnft.pdf}
    % \vspace{-15pt}
    \caption{More quantitative comparisons with DiffusionNFT~\cite{zheng2025diffusionnft}.
    }
    \label{fig: diffusion}
    % \vspace{-0.4cm}
\end{figure*}

\subsection{Failure Cases and Limitations}
\label{sec: limitation}
Despite the superior performance of Flow-OPD across both subjective and objective benchmarks, certain limitations persist. A primary constraint is the \textbf{performance ceiling imposed by teacher models}. As illustrated in Fig.~\ref{fig: failure}, when specialized teachers fail to synthesize semantically correct images, these inaccuracies are propagated through the dense supervisory signals. Such erroneous guidance introduces noise into the distillation objective, ultimately hindering the student’s ability to transcend the inherent limitations of the teacher ensemble.
%
Another inherent limitation is the requirement for \textbf{architectural homogeneity} between the teacher and student models to facilitate fine-grained, step-wise supervision. Looking forward, we aim to explore the broader potential of Flow-OPD through several promising directions, including: (1) \textbf{Co-evolutionary Distillation}, where teachers and students iteratively refine each other; (2) \textbf{Self-Distillation} mechanisms to boost performance without external teachers; and (3) \textbf{Cross-Vocabulary Distillation} to bridge the gap between heterogeneous model architectures.
\begin{figure*}[!h]
    \centering
    \includegraphics[scale=0.45]{images/failure_case.pdf
    }
    % \vspace{-15pt}
    \caption{More quantitative comparisons with DiffusionNFT~\cite{zheng2025diffusionnft}.
    }
    \label{fig: failure}
    % \vspace{-0.4cm}
\end{figure*}


\section*{Broader Impact}
\label{ethics}
Our work on \textbf{Flow-OPD} introduces a robust framework for multi-task alignment in generative models, carrying both positive societal contributions and potential risks that necessitate careful consideration.

\paragraph{Positive Societal Impacts} 
The primary contribution of this work is the enhancement of \textbf{functional reliability} in AI-generated content. By improving layout coherence and OCR accuracy, Flow-OPD can significantly benefit professional fields such as automated graphic design, educational content creation, and assistive technologies for the visually impaired. Furthermore, our Multi-Teacher paradigm promotes a more balanced optimization objective, which mitigates the "winner-takes-all" bias inherent in single-reward reinforcement learning, potentially leading to more diverse and representative generative systems.

\paragraph{Negative Societal Impacts}
Despite these benefits, the increased proficiency in generating high-quality, instruction-following images could be misused for the creation of sophisticated \textbf{disinformation or deceptive visual content}. Although our model inherits the safety filters of its foundation model, the improved structural realism might be exploited to generate more convincing fake documents or misleading social media assets. To mitigate this, we advocate for the integration of digital watermarking and provenance tracking (e.g., C2PA) in downstream applications. Additionally, like all large-scale generative models, there is a risk that the specialized teachers may harbor latent biases present in their training data, which could be inadvertently distilled into the student model. We encourage the community to employ bias-detection benchmarks alongside our framework to ensure equitable performance across all demographics.